
\section{Related Work}

\textbf{FMRI-Image Reconstruction.} Early approaches leveraged linear models to decode fMRI signals into visual features~\citep{kay2008identifying, takagi2023high}. Later work explores deep generative models: approaches such as~\citep{lin2022mind, ozcelik2022reconstruction, goodfellow2020generative} map fMRI signals into the latent space of GANs for image reconstruction, while \textcolor{black}{BraVL~\citep{du2023decoding} employs a multimodal VAE to jointly model relationships between brain activity and visual–linguistic features for neural decoding}. With advances in vision–language models~\citep{radford2021learning, liang2024survey}, more recent approaches map fMRI signals into CLIP’s image-embedding space~\citep{scotti2024mindeye2, scotti2023reconstructing} and then use diffusion models for image reconstruction~\citep{rombach2022high, xu2023versatile, podell2023sdxl}. \textcolor{black}{Building on this diffusion-based direction, MindDiffuser~\citep{lu2023minddiffuser} introduces a two-stage framework that first decodes semantic information and then aligns structural information with CLIP visual features decoded from fMRI.} Unlike prior work that directly maps fMRI signals to joint text–image spaces~\citep{wang2024decoding, quan2024psychometry}, we compare multiple representation spaces and find that text embeddings from language models~\citep{raffel2020exploring} exhibit the strongest alignment with fMRI signals. This insight motivates our approach of reconstructing images via the embedding space of language models. \textcolor{black}{Cross-subject decoding has also gained attention. MindTuner~\citep{gong2025mindtuner} introduces visual-fingerprint modeling using Skip-LoRA to improve cross-subject alignment and semantic correction.
Psychometry~\citep{quan2024psychometry} proposes an omni Mixture-of-Experts architecture that captures both inter-subject shared structure and subject-specific variability, combined with retrieval-enhanced inference for improved reconstruction.
Wills Aligner~\citep{bao2025wills} further advances multi-subject collaborative decoding by combining anatomical alignment with subject-guided Mixture-of-Brain-Expert adapters. Beyond static images, recent work has extended brain decoding to dynamic visual stimuli.
NeuroClips~\citep{gong2024neuroclips} reconstructs such stimuli from fMRI by combining decoded keyframes with low-level perceptual flows to improve temporal consistency.
BrainNEDS~\citep{yeung2025reanimating} utilizes video diffusion models to disentangle static and dynamic components, offering insights into which visual features are prioritized by the brain.} 


% \textbf{Diffusion models} have been widely adopted across a range of tasks, including image generation/editing~\citep{wijmans1995solution, gal2022image, song2020denoising}, and text-to-image synthesis~\citep{ruiz2023dreambooth}.  These generative models learn to synthesize data by reversing a predefined noise-injection process. To enhance image quality, several techniques have been proposed. For example, ControlNet \citep{zhang2023adding} specify some high-level features of images for controlling semantic structures and, and GLIGEN \citep{li2023gligen} design position-aware adapters on top of the diffusion models for spatially-conditioned image generation. Another approach is to use training-free methods to steer diffusion models through manipulating latent or cross-attention maps according to spatial or semantic constraints during inferences without additional training \citep{chen2024training, yang2024mastering}.

\textbf{Diffusion models.} Diffusion models have become foundational in generative tasks like image creation and editing~\citep{gal2022image, song2020denoising}, as well as text-to-image synthesis~\citep{ruiz2023dreambooth}. To enhance control over generated content, ControlNet~\citep{zhang2023adding} introduces high-level image features for controlling and GLIGEN~\citep{li2023gligen,zhang2025layercraft} incorporates position-aware adapters for spatial grounded generation. Meanwhile, there are also training-free methods that adjust latent or attention maps during inference to guide outputs without additional training~\citep{chen2024training, yang2024mastering}. In our work, we guide the diffusion process by modifying cross-attention layers during inference to integrate object-level descriptions derived from fMRI data for image reconstruction.

\textbf{Prompt Optimization.} Prompt optimization aims to discover effective textual prompts for LLMs without model fine-tuning. Gradient-based methods~\citep{shin2020autoprompt, shi2022toward, wen2023hard} update prompts using gradients or differentiable embeddings. Gradient-free approaches treat LLMs as black boxes, using heuristic search~\citep{prasad2022grips, pryzant2023automatic}, reinforcement learning~\citep{deng2022rlprompt, zhang2022tempera}, or evolutionary strategies~\citep{zhou2022large, yang2023large, guo2025evopromptconnectingllmsevolutionary}. We designed our gradient-free prompt optimization based on beam search to optimize attribute keywords for black-box vision-language models.


% \textbf{fMRI-Image Reconstruction.} Traditional methods for fMRI image reconstruction include the use of linear models to predict features from fMRI\citep{kay2008identifying, takagi2023high}. In recent years, deep learning methods have been proposed to improve the quality of fMRI image reconstruction. For example, \citep{lin2022mind, ozcelik2022reconstruction} map fMRI to the latent space of generative adversarial networks (GANs) \citep{goodfellow2020generative} to reconstruct images. After that, due to the advances in vision-language models \citep{radford2021learning, liang2024survey}  researchers explored the mapping fMRI to CLIP \citep{radford2021learning} image space with the help of semantic information \citep{scotti2024mindeye2,wang2024decoding}, and reconstruct the image with the help of diffusion models \citep{rombach2022high, xu2023versatile, podell2023sdxl}. \textcolor{teal}{Despite strong performance, prior methods assume CLIP image features align with brain activity without further assessment. In contrast, we compare multiple representation spaces and find that text representation extracted by T5~\citep{raffel2020exploring} align more closely with fMRI signals. We therefore map fMRI to text space instead of image space for reconstruction, enabling more semantically accurate reconstructions.}
% fmri to text then image, give insights

% \textbf{Diffusion models} have been widely adopted across a range of tasks, including image generation/editing~\citep{wijmans1995solution, gal2022image, song2020denoising}, and text-to-image synthesis~\citep{ruiz2023dreambooth}.  These generative models learn to synthesize data by reversing a predefined noise-injection process. To enhance image quality, several techniques have been proposed. For example, ControlNet \citep{zhang2023adding} specify some high-level features of images for controlling semantic structures and, and GLIGEN \citep{li2023gligen} design position-aware adapters on top of the diffusion models for spatially-conditioned image generation. Another approach is to use training-free methods to steer diffusion models through manipulating latent or cross-attention maps according to spatial or semantic constraints during inference stages without additional training \citep{chen2024training, yang2024mastering}. 

% \textbf{Prompt Optimization} methods aim to automatically discover effective textual prompts to maximize large language models' (LLMs) performance without fine-tuning model parameters. For white-box models, gradient-based methods~\citep{shin2020autoprompt,shi2022toward,wen2023hard} iteratively refine prompts using model gradients, either via differentiable relaxations for discrete prompts or continuous embeddings. Despite their efficiency and strong performance, these approaches require gradient access, limiting applicability to black-box or closed API LLMs. In contrast, gradient-free methods optimize prompts without gradient access by treating the language model as a black-box function. These approaches include heuristic search strategies, such as beam search and iterative editing~\citep{prasad2022grips,pryzant2023automatic}, reinforcement learning methods that use sequential decision-making guided by feedback signals~\citep{deng2022rlprompt,zhang2022tempera}, and evolutionary-inspired algorithms leveraging population-based mutation and selection processes~\citep{zhou2022large,yang2023large,guo2025evopromptconnectingllmsevolutionary}. In our work, we designed our gradient-free prompt optimization based on beam search to find the optimal attribute keyword in the prompt for the black-box VLM.

\begin{table}[t]
\centering
\scriptsize
\caption{\small Top-5 keywords scored by $1-\text{LPIPS}$ before searching and after 10, 20, 30 search steps. The top-5 results remain unchanged after 30 search rounds. The search results indicate a clear preference for keywords related to spatial and positional relations, with most of the top-performing keywords in the final results containing the term '\textit{spatial}'.}
\resizebox{\linewidth}{!}{
\begin{tabular}{c|cccc}
\toprule
\normalsize{\textbf{Rank}} & \normalsize{\textbf{Initial}} & \normalsize{\textbf{Round 10}} & \normalsize{\textbf{Round 20}} & \normalsize{\textbf{Round 30}+}\\ \hline
\normalsize{\#1} & \texttt{Spatial Configuration} & \texttt{Spatial Arrangement} & \texttt{Spatial Organization} & \texttt{Spatial Layout} \\
\normalsize{\#2} & \texttt{Positional Relation} & \texttt{Spatial Configuration} & \texttt{Spatial Structure} & \texttt{Spatial Patterns} \\
\normalsize{\#3} & \texttt{Location Relation} & \texttt{Spatial Interaction} & \texttt{Spatial Arrangement} & \texttt{Spatial Organization} \\
\normalsize{\#4} & \texttt{Descriptive Attribute} & \texttt{Positional Relation} & \texttt{Spatial Configuration} & \texttt{Relative Position} \\
\normalsize{\#5} & \texttt{Inclusion Dependency} & \texttt{Feature Relation} & \texttt{Spatial Interaction} & \texttt{Spatial Relationships} \\
\bottomrule
\end{tabular}
}
\label{tab:search_words}
\end{table}
